\section{Префикс-функция}
Для поиска подстрок и анализа периодичности возвращает массив $\pi$: на позиции $i$ --- длина наибольшего собственного префикса строки $s[0\ldots i]$, совпадающего с её суффиксом; $\pi[0]=0$.
Вход --- строка из любых байтов, включая пустую, длины $n<\texttt{INT\_MAX}$; индексы с нуля, символы Юникода отдельно не декодируются.
Время $O(n)$, память $O(n)$.
\begin{lstlisting}[name=str_prefix]
vector<int> prefixFunction(const string& text) {
    int length = text.size();
    vector<int> prefix(length);
    for (int index = 1; index < length; ++index) {
        int matched = prefix[index - 1];
        while (matched > 0 && text[index] != text[matched])
            matched = prefix[matched - 1];
        if (text[index] == text[matched]) ++matched;
        prefix[index] = matched;
    }
    return prefix;
}
\end{lstlisting}
% \subsection*{Для чего нужен}
% Для каждой позиции $i$ находит длину наибольшего собственного префикса
% строки $s[0\ldots i]$, совпадающего с её суффиксом. В частности,
% $\pi[0]=0$. Нужна для поиска подстроки и анализа периодичности.
% Время $O(n)$, память $O(n)$, где $n$ --- длина строки в байтах.
% Допустимы пустая строка и любые байты; $n<\texttt{INT\_MAX}$.
% Индексы начинаются с нуля; символы Юникода отдельно не декодируются.
% Код адаптирован из \url{https://ru.algorithmica.org/cs/string-searching/prefix-function/}:
% переименованы переменные, строка передаётся по константной ссылке.
% \subsection*{Реализация}
%
% \subsection*{Возможные улучшения и дополнения}
% Для непустого шаблона $p$ вычислите функцию на $p+\#+t$, где разделитель
% не встречается во входе. Значение $|p|$ на позиции $i$ означает вхождение
% в тексте с началом $i-2|p|$. Если свободного байта нет, используйте
% целочисленный алфавит с отдельным разделителем либо потоковый алгоритм
% Кнута--Морриса--Пратта, хранящий только префикс-функцию шаблона.
% Для непустой строки кандидат на период $k=n-\pi[n-1]$;
% строка состоит из целых повторов блока длины $k$, если $n\bmod k=0$.
% Все границы строки перечисляются переходами $k\gets\pi[k-1]$.

\section{Z-функция}
Для поиска подстрок, границ и периодов возвращает массив $z$: на позиции $i$ --- длина наибольшего общего префикса всей строки и суффикса с позиции $i$; принято $z[0]=0$.
Вход --- строка из любых байтов, включая пустую, длины $n<\texttt{INT\_MAX}$; индексы с нуля, сравниваются байты без декодирования Юникода.
Время $O(n)$, память $O(n)$.
\begin{lstlisting}[name=str_z]
vector<int> zFunction(const string& text) {
    int length = text.size();
    vector<int> z(length);
    int left = 0, right = 0;
    for (int index = 1; index < length; ++index) {
        if (index <= right)
            z[index] = min(right - index + 1, z[index - left]);
        while (index + z[index] < length &&
               text[z[index]] == text[index + z[index]])
            ++z[index];
        if (index + z[index] - 1 > right) {
            left = index;
            right = index + z[index] - 1;
        }
    }
    return z;
}
\end{lstlisting}
% \subsection*{Для чего нужна}
% $z[i]$ --- длина наибольшего общего префикса всей строки и её суффикса,
% начинающегося в позиции $i$. Здесь принято $z[0]=0$.
% Применяется для поиска подстрок, границ и периодов строки.
% Время $O(n)$, память $O(n)$, где $n$ --- длина в байтах.
% Допустимы пустая строка и любые байты; $n<\texttt{INT\_MAX}$.
% Индексы с нуля; сравниваются байты, а не символы Юникода.
% Код адаптирован из \url{https://ru.algorithmica.org/cs/string-searching/z-function/}:
% переименованы переменные, строка передаётся по константной ссылке.
% \subsection*{Реализация}
%
% \subsection*{Возможные улучшения и дополнения}
% Поддерживается отрезок $[left,right]$, совпадающий с префиксом.
% Для поиска непустого $p$ в тексте $t$ используйте $p+\#+t$
% с разделителем вне алфавита входа. При $i>|p|$ условие $z[i]=|p|$
% даёт начало вхождения $i-|p|-1$. Для полного байтового алфавита
% перейдите к целым числам с отдельным разделителем.
% Условие $i+z[i]=n$ выделяет суффикс, равный префиксу.
% Число $k$, $1\le k<n$, является периодом тогда и только тогда,
% когда $k+z[k]=n$; для целых повторов дополнительно нужно $n\bmod k=0$.

\section{Алгоритм Манакера: обе чётности}
Находит радиусы всех максимальных палиндромов строки и возвращает массивы $(odd,even)$ за $O(n)$ времени и $O(n)$ памяти.
Нечётный радиус $odd[i]$ включает центр $i$ и задаёт полуинтервал $[i-odd[i]+1,i+odd[i])$ длины $2odd[i]-1$; чётный $even[i]$ имеет центр между $i-1$ и $i$ и задаёт $[i-even[i],i+even[i])$ длины $2even[i]$.
Вход --- любые байты, $n<\texttt{INT\_MAX}$, индексы с нуля, сравнение без декодирования Юникода; пустая строка даёт два пустых массива.
\begin{lstlisting}[name=str_manacher]
pair<vector<int>, vector<int>> manacher(const string& text) {
    int length = text.size();
    vector<int> odd(length), even(length);
    int left = 0, right = -1;
    for (int center = 0; center < length; ++center) {
        int radius = 1;
        if (center <= right)
            radius = min(odd[left + (right - center)], right - center + 1);
        while (center - radius >= 0 && center + radius < length &&
               text[center - radius] == text[center + radius])
            ++radius;
        odd[center] = radius;
        if (center + radius - 1 > right) {
            left = center - radius + 1;
            right = center + radius - 1;
        }
    }
    left = 0;
    right = -1;
    for (int center = 0; center < length; ++center) {
        int radius = 0;
        if (center <= right)
            radius = min(even[left + (right - center) + 1],
                         right - center + 1);
        while (center - radius - 1 >= 0 && center + radius < length &&
               text[center - radius - 1] == text[center + radius])
            ++radius;
        even[center] = radius;
        if (center + radius - 1 > right) {
            left = center - radius;
            right = center + radius - 1;
        }
    }
    return {odd, even};
}
\end{lstlisting}
% \subsection*{Для чего нужен}
% Находит радиусы всех максимальных палиндромов за $O(n)$ времени
% и $O(n)$ памяти.
% Нечётный радиус $odd[i]$ включает центр $i$: палиндром занимает
% $[i-odd[i]+1,i+odd[i])$ и имеет длину $2odd[i]-1$.
% Чётный радиус $even[i]$ задаёт центр между $i-1$ и $i$:
% палиндром занимает $[i-even[i],i+even[i])$, длина $2even[i]$.
% Индексы с нуля; пустая строка даёт два пустых массива.
% Любые байты допустимы, сравнение побайтовое, без декодирования Юникода;
% $n<\texttt{INT\_MAX}$.
% Код адаптирован из \url{https://ru.algorithmica.org/cs/string-searching/manacher/}:
% обе чётности объединены в одну функцию, центр чётного палиндрома
% перенесён на правый символ пары; границы проверяются явно.
% \subsection*{Реализация}
%
% \subsection*{Возможные улучшения и дополнения}
% Максимум длин по двум массивам даёт самый длинный палиндром.
% Сумма всех радиусов в типе \texttt{ll} равна числу непустых
% палиндромных подстрок с учётом позиций, а не числу различных строк.
% Проверка непустого полуинтервала $[l,r)$ за $O(1)$:
% пусть $d=r-l$ и $c=l+\lfloor d/2\rfloor$.
% При нечётном $d$ нужно $odd[c]\ge\lfloor d/2\rfloor+1$,
% при чётном --- $even[c]\ge d/2$.
% Для различных палиндромов и добавления символов онлайн подходит
% палиндромное дерево.

\section{Ахо--Корасик: число вхождений каждого шаблона}
Считает вхождения набора шаблонов в текст, включая перекрытия, и возвращает \texttt{vector<ll>} в порядке шаблонов; дубликаты получают одинаковые ответы, пустой шаблон --- $n+1$, пустые текст и набор допустимы.
Алфавит строго \texttt{a}--\texttt{z}, некорректный байт вызывает \texttt{invalid\_argument}; для суммы длин шаблонов $L$, их числа $k$ и длины текста $n$ требуется $L+1,k,n<\texttt{INT\_MAX}$.
Автомат строится заново для каждого текста; при числе вершин бора $V\le L+1$ время $O(L+26V+n+k)$, дополнительная память $O(26V+k)$.
\begin{lstlisting}[name=str_aho]
vector<ll> ahoCounts(const vector<string>& patterns, const string& text) {
    array<int, 26> empty;
    empty.fill(-1);
    vector<array<int, 26>> next(1, empty);
    vector<int> fail(1);
    vector<ll> count(1);
    vector<int> terminal;
    auto letterIndex = [](char letter) {
        if (letter < 'a' || letter > 'z')
            throw invalid_argument("Ожидается буква a-z");
        return letter - 'a';
    };
    for (const string& pattern : patterns) {
        int vertex = 0;
        for (char letter : pattern) {
            int symbol = letterIndex(letter);
            if (next[vertex][symbol] == -1) {
                int child = (int)next.size();
                next.push_back(empty);
                fail.push_back(0);
                count.push_back(0);
                next[vertex][symbol] = child;
            }
            vertex = next[vertex][symbol];
        }
        terminal.push_back(vertex);
    }
    vector<int> order = {0};
    for (int index = 0; index < (int)order.size(); ++index) {
        int vertex = order[index];
        for (int symbol = 0; symbol < 26; ++symbol) {
            int child = next[vertex][symbol];
            int fallback = vertex == 0 ? 0 : next[fail[vertex]][symbol];
            if (child == -1) {
                next[vertex][symbol] = fallback;
            } else {
                fail[child] = fallback;
                order.push_back(child);
            }
        }
    }
    int vertex = 0;
    for (char letter : text) {
        vertex = next[vertex][letterIndex(letter)];
        ++count[vertex];
    }
    for (int index = (int)order.size() - 1; index > 0; --index) {
        int current = order[index];
        count[fail[current]] += count[current];
    }
    vector<ll> counts;
    for (int endpoint : terminal)
        counts.push_back(endpoint == 0 ? (ll)text.size() + 1 :
                         count[endpoint]);
    return counts;
}
\end{lstlisting}
% \subsection*{Для чего нужен}
% Считает все вхождения набора шаблонов в один текст, включая перекрытия.
%
% Ответ типа \texttt{vector<ll>} идёт в порядке шаблонов;
% дубликаты получают одинаковые ответы. Пустой шаблон встречается $n+1$
% раз, пустые текст и набор шаблонов допустимы.
% Алфавит строго \texttt{a}--\texttt{z}; некорректный байт вызывает
% \texttt{invalid\_argument}. Пусть $L$ --- сумма длин шаблонов,
% $k$ --- их число, $n$ --- длина текста, $V\le L+1$ --- число вершин бора.
% Требуется $L+1,k,n<\texttt{INT\_MAX}$.
% Время $O(L+26V+n+k)$, дополнительная память $O(26V+k)$.
% Автомат строится заново для каждого текста.
% Тематический источник: \url{https://ru.algorithmica.org/cs/string-structures/aho-corasick/}.
% Код ниже самостоятельный: вместо указателей и ленивых переходов из статьи
% использованы массивы, обход в ширину и накопление счётчиков по суффиксным ссылкам.
% \subsection*{Реализация}
%
% \subsection*{Возможные улучшения и дополнения}
% Для пустого шаблона ответ отдельно равен $n+1$: это все позиции между
% символами и по краям текста. Для непустых шаблонов обратный порядок
% обхода в ширину передаёт счётчик вершины её суффиксной ссылке:
% так учитываются все шаблоны-суффиксы без перебора совпадений.
% Для большого алфавита можно хранить разреженные переходы, пересмотрев
% оценки времени. Для выдачи позиций вхождений добавьте выходные ссылки
% на терминальные вершины; время тогда зависит и от числа совпадений.

